% ---------------------------------------------------------------------
% --- Arquivo principal e os demais serao os dos capitulos.
% --- EXPRESSÔES ENTRE <> DEVERÃO SER COMPLETADAS COM A INFORMAÇÂO ESPECÍFICA DO TRABALHO 
% ---------------------------------------------------------------------

\documentclass[ruledheader]{abnt_UFF}

%==\usepackage[colorlinks]{hyperref}
%\usepackage[printonlyused,withpage]{acronym}

%---pacotes para hiphenizacao e acentuacao em portugues
\usepackage[utf8]{inputenc}
\usepackage[brazil]{babel}
\usepackage[T1]{fontenc}


%---fornece a capacidade de criar hiperlinks no documento
%\usepackage[dvips,colorlinks=False, pdfstartview=FitV,citecolor=green,urlcolor=black,plainpages=false, pdftitle={Dissertação de Mestrado Yona Lopes}]{hyperref}%backref,
\usepackage[hidelinks, breaklinks, colorlinks=false, pdfstartview=FitV, linkcolor=black, citecolor=green,urlcolor=black,plainpages=false]{hyperref} % usar  hidelinks para esconder os links, continuara o link mas sem aparecer cor ou borda

\usepackage[breaklinks]{hyperref} % Problema quando o nome da figura/tabela/seção eh grande
%\usepackage[dvips,pdfstartview=FitV]{hyperref}%backref,
%\usepackage[dvips]{color}
	
%---pacotes para criacao de index
\usepackage{makeidx}
%\usepackage[pdftex,colorlinks=true, pdfstartview=FitV, linkcolor=black, citecolor=black,urlcolor=black,plainpages=false]{hyperref}
\usepackage{pdfcolmk}

%--- pacote para figuras
\usepackage{epsf}
\usepackage{epsfig,graphicx}
\usepackage{subfigure}


%--- pacote de simbolos
\usepackage{latexsym}
%\usepackage{textcomp}

%--- simbolos matematicos
\usepackage{amsmath}
%\usepackage{amssymb}

%----- códigos fonte
%\usepackage{listings}
%% opções do pacote listings
%\lstset{numbers=left,
%language=python,
%stepnumber=1,
%firstnumber=1,
%numberstyle=\tiny,
%extendedchars=true,
%breaklines=true,
%frame=tb,
%basicstyle=\footnotesize,
%stringstyle=\ttfamily,
%showstringspaces=false
%backgroundcolor=\color{gray}
%}
%


%--- pacote para gerar pseudo-codigo
%\usepackage{algorithm}
%\usepackage{algorithmic}
% \usepackage[linesnumbered,boxed,lined]{algorithm2e} 
\usepackage[algoruled,portuguese,linesnumbered]{algorithm2e} 
%\floatname{algorithm}{Algoritmo}

%--- outros pacotes
\usepackage{url}
\usepackage{longtable}
\usepackage{lscape}

%\usepackage[notbib, notlof, notlot]{tocbibind}

%Tabela Colorida e pacotes de tabela
%\usepackage{colortbl}
\usepackage{array}

%--- Acronyms
%\usepackage{acronym}
\usepackage[printonlyused,withpage]{acronym}

\usepackage{multicol}
\usepackage{multirow}
\usepackage{rotating}

%--- Notes
\usepackage{todonotes}

\hyphenation{
a-de-qua-da-men-te 
di-men-sio-na-men-to 
re-di-re-cio-na
}

%---------usando tipo de fonte padrao
\renewcommand{\ABNTchapterfont}{\bfseries\fontfamily{cmr}\fontseries{b}\selectfont}
\renewcommand{\ABNTsectionfont}{\bfseries\fontfamily{cmr}}

\newcommand{\todot}[1]{\todo[color=green!40,fancyline]{#1}}


\makeindex
% --- -----------------------------------------------------------------
% --- Documento Principal.
% --- -----------------------------------------------------------------
%\usepackage[pdftex]{hyperref}
%\hypersetup{colorlinks, sitecolor=black, pdftex}
\begin{document}

% --- -----------------------------------------------------------------
% --- Titulo, abstract, dedicatorias e agradecimentos.
% --- Indice geral, lista de figuras e tabelas.
% --- -----------------------------------------------------------------

% --- -----------------------------------------------------------------
% --- Elementos usados na Capa e na Folha de Rosto.
% --- EXPRESSES ENTRE <> DEVERO SER COMPLETADAS COM A INFORMAO ESPECFICA DO TRABALHO
% --- E OS SMBOLOS <> DEVEM SER RETIRADOS 
% --- -----------------------------------------------------------------
\autor{NICOLAS DE SOUSA TEODOSIO E VICTOR HUGO NOVAIS RODRIGUES} % deve ser escrito em maiusculo

\titulo{ANÁLISE DE SENTIMENTO E MINERAÇÃO DE OPINIÕES APLICADO NO TWITTER} % deve ser escrito em maiusculo

\instituicao{INSTITUTO INFNET} % deve ser escrito em maiusculo

\orientador{CASSIUS FIGUEIREDO}% deve ser escrito em maiusculo - preencher com o nome do seu orientador

%\coorientador{NOME} %preencher se houver.

\local{RIO DE JANEIRO} % deve ser escrito em maiusculo

\data{2016} % ano da defesa

\comentario{Trabalho de Conclusão de Curso apresentado ao Programa de Graduação em Engenharia da Computação do \mbox{Instituto Infnet} como parte dos requisitos necessários à obtenção do título de \mbox{Bacharel em Engenharia da Computação}.}%preencha com a sua area de concentracao


% --- -----------------------------------------------------------------
% --- Capa. (Capa externa, aquela com as letrinhas douradas)(Obrigatorio)
% --- ----------------------------------------------------------------
\capa

% --- -----------------------------------------------------------------
% --- Folha de rosto. (Obrigatorio)
% --- ----------------------------------------------------------------
\folhaderosto


\pagestyle{ruledheader}
\setcounter{page}{1}
\pagenumbering{roman}

% --- -----------------------------------------------------------------
% --- Ficha Catalográfica. (Obrigatorio)
% --- ----------------------------------------------------------------
%
%\cleardoublepage
%\thispagestyle{empty}
%\vspace*{130mm}
%
%\begin{flushright}
%
%\hspace{8em}
%\fbox{\begin{minipage}{10cm}

%....

%\end{minipage}}
%
%\end{flushright}
%\newpage

% --- -----------------------------------------------------------------
% --- Termo de aprovacao. (Obrigatorio)
% --- ----------------------------------------------------------------


\cleardoublepage
\thispagestyle{empty}

\vspace{-60mm}

\begin{center}
   {\large NICOLAS DE SOUSA TEODOSIO E VICTOR HUGO NOVAIS RODRIGUES}\\
   \vspace{7mm}

  PySent: ANÁLISE DE SENTIMENTO E MINERAÇÃO DE DADOS APLICADO NO TWITTER\\
%   \vspace{10mm}
 \vspace{8mm}  %diminui para ajustar a data que estava pulando para outra folha
\end{center}

\noindent
\begin{flushright}
\begin{minipage}[t]{8cm} 
%\begin{minipage}{\columnwidth}

Trabalho de Conclusão de Curso apresentado ao Programa de Graduação em Engenharia da Computação do \mbox{Instituto Infnet} como parte dos requisitos
necessários à obtenção do título de \mbox{Bacharel em Engenharia da Computação}

\end{minipage}
\end{flushright}
\vspace{ 5mm}  %original 1cm
\noindent
Aprovada em XX agosto de 2016. \\
\begin{flushright}
  \parbox{15cm}
  {
  \begin{center}
  BANCA EXAMINADORA \\
  \vspace{5mm}
  \rule{11cm}{.1mm} \\
	Prof. Cassius Figueiredo, M.Sc. - Orientador \\ 
	Instituto INFNET \\
  \vspace{5mm}
  \rule{11cm}{.1mm} \\
    A definir \\ 
  \vspace{5mm}
  \rule{11cm}{.1mm} \\
    A definir \\ 
  \end{center}
  }
\end{flushright}
\begin{center}
  \vspace{4mm} % 
 Rio de Janeiro \\
  %\vspace{6mm}
  2016
\end{center}

% --- -----------------------------------------------------------------
% --- Dedicatoria.(Opcional)
% --- -----------------------------------------------------------------
\cleardoublepage
\thispagestyle{empty}
\vspace*{200mm}

\begin{flushright}
{\em 
Às nossas famílias e amigos que durante esses cinco anos nos apoiaram nos momentos mais difíceis.
Aos colegas e mestres que nos acompanharam nesta empreitada que no início eram desconhecidos, e hoje amigos que levaremos para o resto da vida.
%Dedicatória(s): À blá blá%Elemento opcional onde o autor presta homenagem ou dedica seu trabalho (ABNT, 2005).
}
\end{flushright}
\newpage


% --- -----------------------------------------------------------------
% --- Agradecimentos.(Opcional)
% --- -----------------------------------------------------------------
\pretextualchapter{Agradecimentos}
\hspace{5mm}
Agradecemos inicialmente ao professor Cassius Figueiredo pelo total suporte nesse trabalho e por seu papel fundamental na nossa trajetória como alunos de Engenharia da Computação. Aos nossos pais, Rosângela de Sousa Teodosio, José Carlos Teodosio da Silva, Deise Luci Gouveia Novais e Paulo Roberto Rodrigues que nos proporcionaram uma base com valores familiares que nos permitiram seguir sempre nossos estudos, independentemente das dificuldades que surgiram pelo caminho.

Aos colegas de classe que fizemos durante estes 5 anos. Esperamos que, de uma forma ou de outra, todos alcancem o que desejam se seguirem o caminho do trabalho duro e da perseverança.

Aos amigos que nos ajudaram profissionalmente e academicamente, servindo de exemplo e dando conselhos que nos tornaram os profissionais que hoje somos: Ezequiel Bertti, João Leite, Thúlio Costa, Felipe Salvini, Dalton Matos, Alan Dieguez, Jesue Sousa e Cesar Frias.

Eu, Nicolas, agradeço à minha esposa Joyce Kelly Dias da Costa e minha irmã Jéssica de Sousa Teodosio por estar ao meu lado durante todo o curso e suportar todos os percalços que enfrentei até aqui.

Eu, Victor Hugo Novais Rodrigues, agradeço ao meu avô e grande exemplo de força, Manuel Ferreira Novais. Aos padrinhos Mauro Sérgio de Almeida Rubianes, Maysa Alexandrino dos Santos Rubianes e Maria Dionísia Pimentel dos Santos e aos tios Manoel de Carvalho Almeida e Vania Luci Gouveia Novais Almeida por todo carinho e apoio na minha criação e educação. À minha namorada Thays Scapin Coca pelo companheirismo e amor que me permitiram ir mais longe do que eu jamais sonhei.

Um agradecimento especial aos professores que cruzaram nossos caminhos durante o curso sempre transmitindo o conhecimento, experiência e sabedoria da maneira mais nobre possível.

Pedimos desculpas àqueles que porventura tenham sido esquecidos. Esta graduação foi fruto do esforço de muitos e não só das duas pessoas que o encabeçam.
% --- -----------------------------------------------------------------
% --- Resumo em portugues.(Obrigatorio)
% --- -----------------------------------------------------------------
\begin{resumo}

Atualmente a internet e micro blogs em geral têm se tornado uma ferramenta de comunicação poderosa entre seus usuários. Bilhões de pessoas compartilham informações e opiniões a todo momento, fazendo desse espaço um ótimo campo de pesquisas comerciais, acadêmicas e sociológicas.  Como esse fenômeno é relativamente recente, as pesquisas destinadas ao tema são consideradas novas, o Twitter, por exemplo foi criado apenas em 2006 e hoje possui 310 milhões de usuários ativos por mês .


Os principais desafios para aplicação da análise de sentimento estão relacionados a linguagens naturais sensíveis ao contexto, que não trazem resultados satisfatórios quando utilizam-se modelos matemáticos muito simples, sendo necessário um grande investimento de tempo em aperfeiçoar os modelos matemáticos disponíveis e adaptá-los à solução em questão.


O objetivo deste trabalho é explorar o potencial existente em pesquisas de opinião,  que podem ser feitas através de análises nas comunicações utilizando a  língua portuguesa nas redes sociais todos os dias, mostrando que o conteúdo escrito pelos usuários tem uma importância tão grande quanto a de números exibidos em redes sociais. A prática desse trabalho demandou uso de algoritmos matemáticos mais elaborados, aquisição de bases de palavras para o uso desse algoritmo, aplicação de um processo produtor-consumidor e acesso ao Twitter para coleta de dados. Foi realizada uma análise de caso sobre o Oscar 2016 para validar a efetividade desse trabalho. 


O algoritmo utilizado foi o \textit{Naive Bayes}, um algoritmo probabilístico, que realiza a análise de sentimento do oscar no momento de picos do evento, um deles sendo a premiação de melhor ator, atriz e filme. Nesses picos foram feitas análises de quantidade  e  taxa de polaridade. Seus resultados atestam a possibilidade da língua portuguesa como objeto de análise para a mineração de dados de forma relevante.


{\hspace{-8mm} \bf{Palavras-chave}}: Análise de sentimento, mídias sociais, twitter, mineração de opiniões, processamento de linguagem natural, linguagens sensíveis a contexto, naive bayes.

%redefine todas siglas para não usado
\acresetall

\end{resumo}

% --- -----------------------------------------------------------------
% --- Resumo em lingua estrangeira.(Obrigatorio)
% --- -----------------------------------------------------------------
\begin{abstract}
	
Today the internet and micro blogs in general have become a powerful communication tool among its users. Billions of people share information and opinions at all times, making this space a great field of commercial, academic and sociological research. As this phenomenon is relatively recent, research aimed at the topic is considered new, Twitter, for example was created only in 2006 and today has 310 million active users per month.


The main challenges for the application of the feeling analysis are related to context-sensitive natural languages, which do not yield satisfactory results when very simple mathematical models are used, requiring a great investment of time in improving the mathematical models available and adapting them to the solution in question.


The objective of this work is to explore the potential of opinion polls, which can be done through analysis in the communications using the Portuguese language in social networks every day, showing that the content written by users is as important as that of numbers displayed on social networks. The practice of this work required the use of more elaborate mathematical algorithms, acquisition of word bases for the use of this algorithm, application of a producer-consumer process and access to Twitter for data collection. A case study on the Oscar 2016 was conducted to validate the effectiveness of this work.


The algorithm used was the Naive Bayes, a probabilistic algorithm that performs the sentiment analysis of oscar at the moment of the peaks of the event, one of them being the best actor, actress and film award. At these peaks were made analyzes of quantity and polarity rate. Their results attest to the possibility of the Portuguese language as an object of analysis for data mining in a relevant way.	


{\hspace{-8mm} \bf{Palavras-chave}}: Sentiment analisys, social midia, twitter, opnion mining, natural language process, Context-sensitive languages, naive bayes.


%redefine todas siglas para não usado
\acresetall

\end{abstract}

% --- -----------------------------------------------------------------
% --- Lista de figuras.(Opcional)
% --- -----------------------------------------------------------------
\cleardoublepage
\listoffigures
\label{listafiguras}
\acresetall

% --- -----------------------------------------------------------------
% --- Lista de tabelas.(Opcional)
% --- -----------------------------------------------------------------
\cleardoublepage
\label{listatabelas}
\listoftables
\cleardoublepage
\acresetall

% --- -----------------------------------------------------------------
% --- Lista de abreviatura.(Opcional)
%Elemento opcional, que consiste na relação alfabética das abreviaturas e siglas utilizadas no texto, seguidas das %palavras ou expressões correspondentes grafadas por extenso. Recomenda-se a elaboração de lista própria para cada %tipo (ABNT, 2005).
% --- ----------------------------------------------------------------
\cleardoublepage
\pretextualchapter{Lista de Abreviaturas e Siglas}
\begin{acronym}
	\acro{API}{Application Program Interface}
	\acro{PLN}{Processamento de Linguagem Natural}
	\acro{IA}{Inteligência Artificial}
	\acro{REST}{\textit{Representational state transfer}}
	\acro{SOAP}{\textit{Simple Object Access Protocol}}
	\acro{JSON}{\textit{Javascript Object Notation}}
	\acro{WSDL}{\textit{Web Services Description Language}}
	\acro{HTTP}{\textit{Hipertext Transfer Protocol}}
	\acro{XML}{\textit{Extensible Markup Language}}
	\acro{FIFO}{\textit{First-In-First-Out}}
	\acro{NLTK}{\textit{\textit{Natural Language Toolkit}}}
	\acro{SQL}{\textit{Structure Query Language}}
	\acro{NoSQL}{\textit{Non Structure Query Language}}
	\acro{RDBMS}{\textit{Relational database management system}} 
	
\end{acronym}

% --- -----------------------------------------------------------------
% --- Sumario.(Obrigatorio)
% --- -----------------------------------------------------------------
\pagestyle{ruledheader}
\tableofcontents
\acresetall



% --- -----------------------------------------------------------------
% --- Insercao dos capitulos.
% --- -----------------------------------------------------------------
\pagestyle{ruledheader}
\setcounter{page}{1}
\pagenumbering{arabic}

\inputencoding{utf8}

\chapter{Introdução} \label{cap:1}

Através do fenômeno da popularização da Internet vivemos hoje um período conhecido como "Era da conhecimento" \cite{lastres1999informaccao}.
Nesse contexto, redes sociais conhecidas, como Facebook e Twitter se tornaram bastante populares por permitirem a seus usuários acesso à um ambiente onde todos possuem voz e vez para se expressar e por consequência, para se informar sobre tudo que acontece no mundo.
Através de \ac{API} disponibilizadas por essas redes sociais, possuímos fácil acesso à um grande volume de opiniões catalogadas - através de \textit{hashtags} - que podem ser utilizadas em pesquisas de opinião sobre um tema ou assunto específico. Tal cenário apresenta-se como uma grande oportunidade de pesquisa em áreas acadêmicas, sociais e comerciais.
Porém, quando o objeto de estudo é a língua portuguesa, nota-se que a mesma carece de trabalhos e implementações na área de mineração de opiniões e análise de sentimento (REFERÊNCIA). Alguns motivos explicam essa carência: poucos investimentos na área de ciência e engenharia da computação em nosso país e a grande dificuldade que a língua portuguesa apresenta ao ser interpretada através de processamento de linguagem natural. \cite{santos2000projecto}


\section{Motivação e Objetivos}\label{sec:1_inicio}
%A motivação deste trabalho é explorar o potencial contido no conteúdo digital gerado todos os dias em redes sociais por %usuários brasileiros. Tais dados possuem informações valiosas que podem ser explorados de inúmeras maneiras



\section{Principais contribuições}\label{sec:1_principais_contribuicoes}

\section{Recursos utilizados}\label{sec:1_recursos_utilizados}

\section{Organização do trabalho}\label{sec:1_org}

Este trabalho está estruturado em 5 capítulos da seguinte forma: no Capítulo~\ref{cap:referencial_teorico}, para embasamento teórico, são apresentados os conceitos de (CONTINUA). Em seguida, no Capítulo~\ref{cap:proposta} , é feita uma análise sobre os principais trabalhos relacionados ao uso dos ... . No Capítulo~\ref{cap:referencial_teorico}, os conceitos do arcabouço utilizado ... , são descritos. Nesse capítulo são mostrados os motivos para a escolha desse arcabouço, .... A proposta XXX é apresentada no Capitulo~\ref{cap:proposta}, onde a arquitetura da proposta é detalhada, assim como seus componentes e algoritmos. Em seguida, o Capítulo~\ref{cap:resultados} apresenta as ferramentas utilizadas para implementação da proposta, o ambiente implementação, a descrição dos experimentos e os principais resultados obtidos com o XXX, assim como a análise dos valores encontrados. Por fim, o Capítulo~\ref{cap:conc} conclui este trabalho, ressaltando os objetivos alcançados com as propostas. As principais vantagens e desvantagens da proposta são discutidas, assim como alguns trabalhos futuros que podem ser desenvolvidos. 

\chapter{Referencial Teórico}\label{cap:XX}

Normalmente o Capítulo~\ref{cap:XX} é parte do referencial teórico/estudo da arte/histórico... Aqui você precisa dar conhecimento para o leitor entender a sua proposta, deixando claro qual o problema. Pode ter mais de um capítulo no caso.


Continue usando citações como esta de exemplo~\cite{fangxing2010,gungor2011}.

Para itens estrangeiros como \textit{feedback}  não esqueça de colocar em itálico. Afirmações como essa: são os mesmos de 40 anos atrás~\cite{gungor2011}, precisam de referencia. 

\section{para seção}\label{sec:XX} %usar o label ajuda muito 

Usar o pacote de acronimos ajuda e muito, pois na primeira vez que o item aparecer no texto ele vai colocar o significado, e depois só a sigla ex: O \ac{SEP} tem como objetivo...

Para enumerar:

\begin{enumerate}
\item Geração: 
na geração, as usinas são classificadas conforme os recursos que utilizam, podendo ser hidroelétricas, termoelétricas, eólicas, nucleares, etc. Nesta fase, certo tipo de energia é transformado em energia elétrica. %A geração de energia inclui as instalações para esta finalidade. 
No caso do uso de hidrelétricas, envolve o armazenamento de um fluido, normalmente água de rios, conversão de energia hidráulica do fluido em energia mecânica em uma turbina hidráulica e a conversão da energia mecânica em energia elétrica por um gerador elétrico. 
\item Transmissão e Distribuição:
condução da energia elétrica das usinas de geração até os consumidores. Após gerada, a tensão é aumentada para que se evite perdas, em seguida, a energia elétrica é transportada em altas tensões, por fim, é distribuída ao consumidor final em baixas tensões. 
\end{enumerate}

Para siglas no plural use o acp: Em resumo, os \acp{TP} e \acp{TC} traduzem. 

Outro exemplo:  \acp{IED}, são capazes não só...

O label nas seções ajuda muito, veja: Conforme descrito na Seção~\ref{sec:XX}, ou no Capítulo~\ref{cap:XX}.

A Figura~\ref{fig:2_juncaomodelodados} precisa sempre vir referenciada.

\begin{figure}[ht]
	\begin{center}
		\centering
		\includegraphics[width=0.5\linewidth]{figuras/images}
	\end{center}	
	\caption{-----}
	\label{fig:2_juncaomodelodados}
\end{figure}

Exemplo de Figura lado a lado na Figura~\ref{fig:2_principioscom}.

\begin{figure*}[ht]
\begin{center}
\mbox{
	\subfigure[\ac{TPAA}, baseado em mensagens \textit{unicast}.]{\epsfig{figure=./figuras/images, width=0.4\textwidth}
	\label{subfig:2_principioscoma}
					}
	\hspace{0.2cm}
	\subfigure[\ac{MCAA}, baseado em mensagens \textit{multicast}.]{\epsfig{figure=./figuras/images, width=0.4\textwidth}
	\label{subfig:2_principioscomb}
					}
	}
\end{center}
\caption{Princípios de comunicação TPAA e MCAA da norma IEC 61850~\cite{fangxing2010}.}
\label{fig:2_principioscom}
\end{figure*}

E um exemplo de tabela, conforme Tabela ~\ref{tab:2_tabmsg}.

\begin{table}[h!]	
	\begin{center}
		\caption{Tipos de mensagens suportadas pelo padrão IEC 61850~\cite{fangxing2010}.}	
		%\setlength{\extrarowheight}{1.5pt}
		\begin{tabular}{|c|p{3.5cm}|p{4cm} |c|c|} \hline
		\label{tab:2_tabmsg}
{\bf Tipo} & {\bf Descrição} & {\bf Exemplo} & {\bf Classes} & {\bf Mensagem e }	\\
		   &			  &				  &	& Requisitos Temporais \\ \hline
		\hline
        1A & Mensagens Rápidas 				& {\it Trips} 					& P2 e P3 & GOOSE (3ms)		\\ \cline{4-5}
           &  				 				& 								& P1 	  & GOOSE (10ms)	\\ \hline 
        1B & Mensagens Rápidas 			 	& Comandos, 					& P2 e P3 & GOOSE (20ms)  	\\ \cline{4-5}
           & (outras)					 	& Mensagens Simples 			& P1	  & GOOSE (100ms) 	\\ \hline
         2 & Velocidade Média 				& Valores de Medidas 			& -  & MMS (100ms)			\\ \hline
         3 & Velocidade Baixa 				& Parâmetros 					& -  & MMS (500ms)			\\ \hline
         4 & Rajada de Dados                & Saída de dados dos 			& P2 e P3 & SV (3ms) 		\\ \cline{4-5}
		       &							&  transformadores				& P1      & SV (10ms)		\\ \hline
         5 & Transferência de Arquivos		& Arquivos grandes 				& -  & MMS($\geq$1000ms)	\\ \hline
        6A & Sincronização de Tempo A 		& Sincronização (\textit{station bus}) 		& -  & TimeSync	\\  \hline
        6B & Sincronização de Tempo A 		& Sincronização (\textit{process bus}) 		& -  & TimeSync \\ \hline
        7 & Mensagem de Comando 			& Comandos da estação HMI 		& -  & MMS(500ms)			\\ \hline
		\end{tabular}
	\end{center}
\end{table}


\chapter{A proposta} \label{cap:prop}

Definição da sua proposta.

Se for apresentar os algoritmos
use por exemplo:

\begin{algorithm} 
	%\begin{algoendfloat} 
	\input{algoritmo_multitreev2.tex} 
	\caption{Algoritmo de cálculo de árvores \textit{multicast} e estabelecimento de rotas} 
	\label{algo:5_multi_tree} 
\end{algorithm}
%\end{algoendfloat}
%\vspace{-0.3cm}

E explique cada linha.


\chapter{Implementação e Análise da proposta}\label{cap:imp}

Descreva toda a sua implementação e depois faça a analise. Dependendo do tamanho podem estar em capitulos separados.

Para colocar apêndice informe em que está no Apêndice~\ref{apend:saidagerador}.

\section{Cenários e parâmetros de teste}\label{simu}


\section{Experimentos realizados e resultados}\label{sec:7_resul}


\chapter{Conclusões} \label{cap:conc}

Um paragrágo relembrando a importancia do cenário

Esse trabalho identificou e abordou alguns desses problemas, assim como propôs, desenvolveu e avaliou um serviço de gerenciamento eXXXXX
Relembrar o que o trabalho fez.


A proposta, XXX,  se destacou pelo XXXX que apresentou quando comparada XXXX. 

A proposta atingiu os seguintes objetivos, exemplo:
\begin{itemize}
\item permitiu que sejam usados IEDs mais simples pois a solução não precisa ser implementada nesses dispositivos;
\item reduziu o tempo de convergência dos algoritmos, o atraso na entrega de dados e o tráfego na rede;
\item atendeu aos requisitos da Norma IEC 61850;
\item implementou e testou um encaminhamento \textit{multicast} independente de camadas e transparente aos dispositivos finais;
\item permitiu uma configuração da rede facilitada;
\item usou o arquivo SCD da norma para autoconfiguração da rede de Telecomunicações;
\item tornou a rede menos sujeita à erros por ser automático;
\item permitiu o uso mais inteligente de recuperação de falhas;
\item permitiu o alcance de tempos de resposta menores por possuir uma característica proativa.
\end{itemize}

Os experimentos e as análises realizadas mostraramXXXXXX

Falar de todos os resultados encontrados de forma sumarizada, máximo de uma folha.

Os testes mostraram, também, que 


Outro ganho relacionado ao uso da técnica....

A análise realizada mostra que ...

\section{Trabalhos Futuros}\label{sec:8_trabfut}

Como trabalhos futuros, pretende-se ...

Uma outra questão é o estudo, desenvolvimento e implementação ...

Por fim, pretende-se fazer ...


% --- -----------------------------------------------------------------
% --- Referencias Bibliograficas. (Obrigatorio)
% --- -----------------------------------------------------------------
\cleardoublepage
%\bibliographystyle{acm-2} 
%\bibliographystyle{abnt-num} % abbrv - abnt-num
\inputencoding{latin1}
\bibliographystyle{IEEEtran}
%\bibliographystyle{uff-ic}
\bibliography{referencia} % arquivo fonte com a bibliografia
%\bibliography{articles,books,standards,websites} % arquivo fonte com a bibliografia


% --- -----------------------------------------------------------------
% --- Apendice.(Opcional)
% --- -----------------------------------------------------------------
\cleardoublepage
\appendix
\inputencoding{latin1}
\chapter{Exemplo da Sa�da do Gerador GOOSE}
\label{apend:saidagerador}

Coloque aqui se apendice

%\addcontentsline{toc}{part}{Index}
%\printindex
\label{LastPage}
\end{document}